\section*{3.2 Stieltjes Measure Function}

We begin by considering a probability space $(\mathbb{R}, \mathcal{B}(\mathbb{R}), P)$ defined on the real number line $\mathbb{R}$, equipped with the Borel algebra $\mathcal{B}(\mathbb{R})$. The distribution function provides a convenient way to visualize a probability measure.

\begin{defbox}{3.4}
Let $(\mathbb{R}, \mathcal{B}(\mathbb{R}), P)$ be a probability space. The \textit{distribution function} induced by $P$ is defined as
\[
F(x) \triangleq P((-\infty, x]) \qquad \text{for } x \in \mathbb{R}.
\]
\end{defbox}


\textbf{Distribution Function Vs. Cumulative Distribution Function} The function $F(x)$ in Definition 3.4 is very similar to the cumulative distribution function (cdf) of a random variable. The main difference is that a cdf is defined for a random variable, whereas a distribution function is defined for a probability measure.

\textbf{Convention} \\
Some people prefer to define the distribution function using open intervals of the form $(-\infty, x)$, instead of the left-open-right-closed intervals $(-\infty, x]$. In this case, the distribution function is defined as $P((-\infty, x))$ for $x \in \mathbb{R}$. The choice between using one type of intervals or the other is a matter of convention.

\begin{thmbox}{3.2}
The distribution function $F(x)$ of a probability measure satisfies the following properties:
\begin{enumerate}
    \item $F(x)$ is non-decreasing.
    \item $F(x)$ is continuous from the right.
    \item $\lim_{x \to \infty} F(x) = 1$.
    \item $\lim_{x \to -\infty} F(x) = 0$.
\end{enumerate}
\end{thmbox}

\noindent\textit{Proof} Property 1 follows from the monotonic property of measure. Property 2 is a consequence of the upper semi-continuity. If $(x_i)_{i=1}^{\infty}$ is a decreasing sequence converging to $x$ from the right, then $(-\infty, x_i]$ is a decreasing sequence of sets, and the intersection is $(-\infty, x]$. Hence, $\lim_{i \to \infty} F(x_i) = \lim_{i \to \infty} P((-\infty, x_i]) = P((-\infty, x]) = F(x)$. Properties 3 and 4 also follow from semi-continuity properties of probability measure. \hfill $\blacksquare$

Motivated by the previous theorem, we define a Stieltjes measure function in general.

\begin{defbox}{3.5}
A function $F: \mathbb{R} \to \mathbb{R}$ is called a \textit{Stieltjes measure function} if:
\begin{enumerate}
    \item $F$ is non-decreasing.
    \item $F$ is continuous from the right, i.e., $\lim_{y \to x^+} F(y) = F(x)$ for all $x \in \mathbb{R}$.
\end{enumerate}
\end{defbox}
A Stieltjes measure function does not necessarily have a density function on the real number line. However, if $F$ is both continuous and differentiable, then the derivative $f(x) = F'(x)$ is a density function, and an interval $[a, b]$ in the real number line has measure $\int_{a}^{b} f(x) \, dx$.

In general, as we will show in the next section, a Stieltjes measure function defines a measure on the Borel algebra of $\mathbb{R}$, called the Lebesgue--Stieltjes measure, which assigns to each left-open-right-closed interval $(a, b]$ the value $F(b) - F(a)$.

\textbf{Example 3.2.1 (Examples of Stieltjes Measure Functions)}
\begin{enumerate}
    \item The identity function $F(x) = x$ is a Stieltjes measure function, since it is both continuous and non-decreasing.
    \item Fix a constant $x_0$. The function
    \[
    F(x) = \begin{cases} 1 & \text{if } x \ge x_0 \\ 0 & \text{if } x < x_0 \end{cases}
    \]
    defines a Stieltjes measure function. The function $F(x)$ is continuous everywhere except at $x = x_0$, where it has a discontinuity jump.
    \item The cumulative distribution function $F(x)$ of the standard normal distribution
    \[
    F(x) = \int_{-\infty}^{x} \frac{1}{\sqrt{2\pi}} e^{-t^2/2} \, dt
    \]
    is a Stieltjes measure function. In fact, if $f(x)$ is the pdf of any continuous-type distribution, the function $F(x) = \int_{-\infty}^{x} f(t) \, dt$ is a Stieltjes measure function.
\end{enumerate}